\documentclass[aspectratio=169]{beamer}

\usetheme{Madrid}
\usepackage[spanish]{babel}
\usepackage[utf8]{inputenc}
\usepackage{textcomp}
\usepackage[T1]{fontenc}
\usepackage{lmodern}
\usepackage{graphicx}
\usepackage{booktabs}
\usepackage{verbatim}
\usepackage{amssymb}
\usepackage{pifont}
\usepackage{moreverb}
\usepackage[most]{tcolorbox}
\usepackage{tikz}
\usetikzlibrary{shapes, arrows, positioning}

% Configuración de TikZ
\tikzstyle{block} = [
  rectangle, draw, rounded corners,
  text centered, minimum height=1cm,
  minimum width=2.5cm, font=\small, fill=gray!5
]
\tikzstyle{arrow} = [->, thick]

% Configuración de tcolorbox
\tcbset{
  colback=gray!5, colframe=gray!50,
  boxrule=0.4pt, sharp corners,
}

\title[BookMate]{BookMate — Sistema de Recomendación de Libros Académicos}
\subtitle{Demo Final, Planes de Pruebas, Despliegue y Configuración}
\author[GRUPO 6]{Delgado R., G. \and Osorio M., A. \and Rojas A., J. \and Torres R., J. \and Valverde G., Y. \and Villanueva A., F.}
\institute[FC-UNI]{UNIVERSIDAD NACIONAL DE INGENIERÍA}
\titlegraphic{\includegraphics[height=1.4cm]{Imagenes/logo_uni.png}}
% NOTA: Reemplazar 'example-image' con 'Imagenes/logo_uni.png' cuando esté disponible

\date{Noviembre 2025}

\begin{document}

% Portada
\begin{frame}
\titlepage
\end{frame}

% Agenda
\begin{frame}
\frametitle{Agenda}
\tableofcontents
\end{frame}

\section{Demo del Sistema}

\begin{frame}
\frametitle{Guía de Demostración}
\begin{columns}
  \begin{column}{0.5\textwidth}
    \textbf{Objetivo:}
    \begin{itemize}
      \item Demostrar funcionalidades principales
      \item Validar casos de uso críticos
      \item Mostrar sistema de recomendaciones IA
      \item Verificar integración de componentes
    \end{itemize}
  \end{column}
  \begin{column}{0.5\textwidth}
    \begin{center}
      \includegraphics[width=0.9\textwidth]{Imagenes/demo_flujo.png}
      \captionof{figure}{Flujo de Demostración (Ejemplo)}
      \vspace{0.2cm}
      \small \textit{Descripción: Diagrama que muestra el flujo de la demostración, desde el inicio de sesión hasta la obtención de recomendaciones con IA.}
    \end{center}
  \end{column}
\end{columns}
\end{frame}

\begin{frame}
\frametitle{Escenarios de Demostración}
\begin{tcolorbox}[title=Escenario 1: Búsqueda y Recomendaciones IA]
\textbf{Flujo:}
\begin{enumerate}
  \item Usuario busca libro "Inteligencia Artificial"
  \item Selecciona libro del catálogo
  \item Solicita recomendaciones (IA)
  \item Sistema genera 6 recomendaciones semánticas
  \item Usuario revisa justificaciones
\end{enumerate}
\end{tcolorbox}

\begin{tcolorbox}[title=Escenario 2: Fallback a Heurísticas]
\textbf{Flujo:}
\begin{enumerate}
  \item Servicio IA no disponible
  \item Sistema detecta fallo automáticamente
  \item Genera recomendaciones heurísticas
  \item Usuario recibe recomendaciones basadas en género/autor
\end{enumerate}
\end{tcolorbox}
\end{frame}

\begin{frame}
\frametitle{Flujos Críticos}
\begin{center}
  \includegraphics[width=0.9\textwidth]{Imagenes/flujo_critico_uc06.png}
  \captionof{figure}{Flujo Crítico: Obtener Recomendaciones con IA (Ejemplo)}
  \vspace{0.2cm}
  \small \textit{Descripción: Diagrama de flujo que muestra el proceso completo de obtención de recomendaciones, incluyendo la verificación de disponibilidad del servicio IA, el cálculo de similitud semántica, y el fallback a heurísticas si es necesario.}
\end{center}
\end{frame}

\section{Plan de Pruebas}

\begin{frame}
\frametitle{Estrategia de Pruebas}
\begin{columns}
  \begin{column}{0.5\textwidth}
    \textbf{Niveles:}
    \begin{itemize}
      \item \textbf{Unitarias:} Clases y métodos individuales
      \item \textbf{Integración:} Interacción entre componentes
      \item \textbf{Sistema:} End-to-end
      \item \textbf{Rendimiento:} Tiempos de respuesta
      \item \textbf{Seguridad:} Autenticación y autorización
    \end{itemize}
  \end{column}
  \begin{column}{0.5\textwidth}
    \begin{center}
      \includegraphics[width=0.9\textwidth]{Imagenes/piramide_pruebas.png}
      \captionof{figure}{Pirámide de Pruebas (Ejemplo)}
      \vspace{0.2cm}
      \small \textit{Descripción: Pirámide que muestra la distribución de pruebas (más unitarias en la base, menos pruebas de sistema en la cima).}
    \end{center}
  \end{column}
\end{columns}
\end{frame}

\begin{frame}
\frametitle{Cobertura de Pruebas}
\begin{table}
\centering
\small
\begin{tabular}{lcc}
\toprule
\textbf{Componente} & \textbf{Cobertura} & \textbf{Objetivo} \\
\midrule
LibrosService & >90\% & \checkmark \\
RecommendationService & >85\% & \checkmark \\
AIServiceClient & >80\% & \checkmark \\
HeuristicRecommender & >85\% & \checkmark \\
Servicio IA (Python) & >80\% & \checkmark \\
\textbf{Total} & \textbf{>80\%} & \checkmark \\
\bottomrule
\end{tabular}
\end{table}

\vspace{0.3cm}

\begin{tcolorbox}[title=Herramientas]
\begin{itemize}
  \item \textbf{Java:} JUnit 5, JaCoCo
  \item \textbf{Python:} pytest, pytest-cov
  \item \textbf{Integración:} Spring Boot Test, Testcontainers
\end{itemize}
\end{tcolorbox}
\end{frame}

\begin{frame}
\frametitle{Pruebas del Servicio IA}
\begin{center}
  \includegraphics[width=0.9\textwidth]{Imagenes/pruebas_ia.png}
  \captionof{figure}{Estructura de Pruebas del Servicio IA (Ejemplo)}
  \vspace{0.2cm}
  \small \textit{Descripción: Diagrama que muestra las pruebas del servicio IA, incluyendo pruebas de generación de embeddings, cálculo de similitud, y endpoints REST.}
\end{center}

\vspace{0.3cm}

\begin{itemize}
  \item Pruebas de generación de embeddings
  \item Pruebas de cálculo de similitud semántica
  \item Pruebas de endpoints REST
  \item Pruebas de integración con base de datos
\end{itemize}
\end{frame}

\section{Plan de Despliegue}

\begin{frame}
\frametitle{Arquitectura de Despliegue con Docker}
\begin{center}
  \includegraphics[width=0.9\textwidth]{Imagenes/despliegue_docker.png}
  \captionof{figure}{Arquitectura de Despliegue Docker Compose (Ejemplo)}
  \vspace{0.2cm}
  \small \textit{Descripción: Diagrama que muestra la arquitectura de despliegue con Docker Compose, incluyendo los contenedores (Backend Spring Boot, Servicio IA Python, PostgreSQL), la red Docker, y los volúmenes de persistencia.}
\end{center}
\end{frame}

\begin{frame}
\frametitle{Configuración Docker Compose}
\begin{columns}
  \begin{column}{0.5\textwidth}
    \textbf{Servicios:}
    \begin{itemize}
      \item \textbf{postgres:} Base de datos PostgreSQL 16
      \item \textbf{backend:} Aplicación Spring Boot
      \item \textbf{ai-service:} Microservicio IA (Python)
    \end{itemize}
    
    \vspace{0.3cm}
    
    \textbf{Características:}
    \begin{itemize}
      \item Health checks configurados
      \item Dependencias entre servicios
      \item Volúmenes para persistencia
      \item Red Docker aislada
    \end{itemize}
  \end{column}
  \begin{column}{0.5\textwidth}
    \begin{center}
      \includegraphics[width=0.9\textwidth]{Imagenes/docker_compose.png}
      \captionof{figure}{Configuración Docker Compose (Ejemplo)}
      \vspace{0.2cm}
      \small \textit{Descripción: Captura de pantalla o diagrama que muestra la estructura del archivo docker-compose.yml con los servicios, variables de entorno, y configuraciones.}
    \end{center}
  \end{column}
\end{columns}
\end{frame}

\begin{frame}
\frametitle{Migraciones Flyway}
\begin{center}
  \includegraphics[width=0.9\textwidth]{Imagenes/flyway_migrations.png}
  \captionof{figure}{Estructura de Migraciones Flyway (Ejemplo)}
  \vspace{0.2cm}
  \small \textit{Descripción: Diagrama que muestra la estructura de migraciones Flyway (V1\_\_Initial\_schema.sql, V2\_\_Add\_authors\_table.sql, etc.) y cómo se ejecutan automáticamente al iniciar la aplicación.}
\end{center}

\vspace{0.3cm}

\begin{itemize}
  \item Versionado automático de esquema
  \item Migraciones ejecutadas al inicio
  \item Historial de cambios en BD
  \item Rollback controlado
\end{itemize}
\end{frame}

\begin{frame}
\frametitle{Estrategia de Backup}
\begin{columns}
  \begin{column}{0.5\textwidth}
    \textbf{Frecuencia:}
    \begin{itemize}
      \item Backups diarios automáticos
      \item Retención: 7 días
      \item Comprimidos (gzip)
      \item Almacenados en volumen persistente
    \end{itemize}
    
    \vspace{0.3cm}
    
    \textbf{Proceso:}
    \begin{enumerate}
      \item Ejecutar \texttt{pg\_dump}
      \item Comprimir backup
      \item Eliminar backups antiguos
    \end{enumerate}
  \end{column}
  \begin{column}{0.5\textwidth}
    \begin{center}
      \includegraphics[width=0.9\textwidth]{Imagenes/backup_strategy.png}
      \captionof{figure}{Estrategia de Backup (Ejemplo)}
      \vspace{0.2cm}
      \small \textit{Descripción: Diagrama que muestra el flujo de backup automático, incluyendo la ejecución diaria, compresión, y eliminación de backups antiguos.}
    \end{center}
  \end{column}
\end{columns}
\end{frame}

\section{Plan de Administración de Configuración}

\begin{frame}
\frametitle{Git Workflow}
\begin{center}
  \includegraphics[width=0.9\textwidth]{Imagenes/git_workflow.png}
  \captionof{figure}{Git Workflow (Ejemplo)}
  \vspace{0.2cm}
  \small \textit{Descripción: Diagrama que muestra el flujo de Git (feature branches → develop → release → main), incluyendo las estrategias de merge y las reglas de cada rama.}
\end{center}

\vspace{0.3cm}

\begin{itemize}
  \item \textbf{main:} Código estable para producción
  \item \textbf{develop:} Integración de features
  \item \textbf{feature/*:} Desarrollo de funcionalidades
  \item \textbf{hotfix/*:} Correcciones urgentes
\end{itemize}
\end{frame}

\begin{frame}
\frametitle{Versionado Semántico}
\begin{table}
\centering
\small
\begin{tabular}{lcc}
\toprule
\textbf{Versión} & \textbf{Tipo de Cambio} & \textbf{Ejemplo} \\
\midrule
\textbf{MAJOR} & Cambios incompatibles & 2.0.0 \\
\textbf{MINOR} & Nuevas funcionalidades & 1.1.0 \\
\textbf{PATCH} & Correcciones de bugs & 1.1.1 \\
\bottomrule
\end{tabular}
\end{table}

\vspace{0.3cm}

\begin{tcolorbox}[title=Ejemplos]
\begin{itemize}
  \item \texttt{v1.0.0}: Versión inicial
  \item \texttt{v1.1.0}: Agregar recomendaciones IA
  \item \texttt{v1.1.1}: Corregir bug en cálculo de similitud
  \item \texttt{v2.0.0}: Refactorización mayor
\end{itemize}
\end{tcolorbox}
\end{frame}

\begin{frame}
\frametitle{Gestión de Dependencias}
\begin{columns}
  \begin{column}{0.5\textwidth}
    \textbf{Java (Maven):}
    \begin{itemize}
      \item Versiones en \texttt{pom.xml}
      \item Actualización con \texttt{mvnw versions:display-dependency-updates}
      \item Gestión centralizada
    \end{itemize}
    
    \vspace{0.3cm}
    
    \textbf{Python:}
    \begin{itemize}
      \item \texttt{requirements.txt}
      \item \texttt{pip freeze > requirements.txt}
      \item Versionado explícito
    \end{itemize}
  \end{column}
  \begin{column}{0.5\textwidth}
    \begin{center}
      \includegraphics[width=0.9\textwidth]{Imagenes/dependencias.png}
      \captionof{figure}{Gestión de Dependencias (Ejemplo)}
      \vspace{0.2cm}
      \small \textit{Descripción: Diagrama que muestra la estructura de dependencias del proyecto, incluyendo las dependencias de Java (Spring Boot, JPA, etc.) y Python (Flask, Sentence Transformers, etc.).}
    \end{center}
  \end{column}
\end{columns}
\end{frame}

\begin{frame}
\frametitle{CI/CD (Opcional)}
\begin{center}
  \includegraphics[width=0.9\textwidth]{Imagenes/cicd_pipeline.png}
  \captionof{figure}{Pipeline CI/CD (Ejemplo)}
  \vspace{0.2cm}
  \small \textit{Descripción: Diagrama que muestra el pipeline CI/CD (Build → Test → Lint → Deploy), incluyendo las herramientas utilizadas (GitHub Actions, Maven, pytest).}
\end{center}

\vspace{0.3cm}

\begin{itemize}
  \item \textbf{Build:} Compilación del código
  \item \textbf{Test:} Ejecución de pruebas
  \item \textbf{Lint:} Verificación de código
  \item \textbf{Deploy:} Despliegue automático (opcional)
\end{itemize}
\end{frame}

\section{Resumen y Conclusiones}

\begin{frame}
\frametitle{Resumen de la Semana 16}
\begin{columns}
  \begin{column}{0.5\textwidth}
    \textbf{Entregables:}
    \begin{itemize}
      \item \checkmark Guía de demostración
      \item \checkmark Plan de pruebas (cobertura >80\%)
      \item \checkmark Plan de despliegue (Docker Compose)
      \item \checkmark Plan de administración de configuración
    \end{itemize}
  \end{column}
  \begin{column}{0.5\textwidth}
    \begin{center}
      \includegraphics[width=0.8\textwidth]{Imagenes/resumen_semana_16.png}
      \captionof{figure}{Icono de Resumen (Ejemplo)}
      \vspace{0.2cm}
      \small \textit{Descripción: Un icono o gráfico que simboliza la finalización de la fase de pruebas, despliegue y configuración del proyecto.}
    \end{center}
  \end{column}
\end{columns}

\vspace{0.3cm}

\begin{tcolorbox}[title=Logros Principales]
\begin{itemize}
  \item Sistema probado con cobertura >80\%
  \item Despliegue reproducible con Docker Compose
  \item Estrategia de backup implementada
  \item Gestión de configuración organizada
\end{itemize}
\end{tcolorbox}
\end{frame}

\begin{frame}
\frametitle{Próximos Pasos}
\begin{itemize}
  \item \textbf{Demo Final:} Presentación del sistema funcional
  \item \textbf{Link al Drive:} Documentación completa y software
  \item \textbf{Indicaciones de Ejecución:} Guía para usuarios finales
  \item \textbf{Subsanación (S17):} Correcciones si aplican
\end{itemize}

\vspace{0.5cm}

\begin{tcolorbox}[title=Estado del Proyecto]
El sistema BookMate está listo para demostración, con:
\begin{itemize}
  \item Funcionalidades principales implementadas
  \item Sistema de recomendaciones IA operativo
  \item Pruebas con cobertura adecuada
  \item Despliegue automatizado
\end{itemize}
\end{tcolorbox}
\end{frame}

\begin{frame}
\frametitle{¡Preguntas?}
\begin{center}
  \Large ¡Gracias por su atención! \\
  \vspace{1cm}
  \textbf{¿Preguntas sobre la demo, pruebas, despliegue o configuración?}
  \vspace{1cm}
  \normalsize
  \textbf{Grupo 6}\\
  Universidad Nacional de Ingeniería\\
  CC341 - Ingeniería de Software\\
  Ciclo 2025-2
  \vspace{0.5cm}
  \textbf{Próxima entrega:} Semana 17 - Subsanación (si aplica)
\end{center}
\end{frame}

\end{document}

